% Bài làm buổi 12/11
\begin{problem}[8DS-HSBN-001]
Thay từng giá trị của $x$ vào công thức $f(x) = 3x - 2$.

$f(0) = 3 \cdot 0 - 2 = -2$.

$f(2) = 3 \cdot 2 - 2 = 6 - 2 = 4$.

$f\left(-\dfrac{1}{3}\right) = 3 \cdot \left(-\dfrac{1}{3}\right) - 2 = -1 - 2 = -3$.

Đáp số: $f(0) = -2$, $f(2) = 4$, $f\left(-\dfrac{1}{3}\right) = -3$
\end{problem}

\begin{problem}[8DS-HDT-013]
a) Diện tích mảnh đất cũ là $a^2$ (m$^2$). Mảnh đất mới có chiều dài $a + 4$ (m), chiều rộng $a - 4$ (m) nên có diện tích $(a + 4)(a - 4) = a^2 - 16$ (m$^2$).

Hiệu hai diện tích: $(a^2 - 16) - a^2 = -16 < 0$. Vậy mảnh đất mới nhỏ hơn mảnh đất cũ $16$ m$^2$.

b) Với $a = 54$: chiều dài $58$, chiều rộng $50$, diện tích là $58 \cdot 50 = 2\,900$.

Kiểm tra: $54^2 - 16 = 2\,916 - 16 = 2\,900$ (khớp).

Đáp số: a) nhỏ hơn $16$ m$^2$; b) $2\,900$
\end{problem}
